\documentclass[12pt,a4paper]{article}
\usepackage{amsmath}
\usepackage{geometry}
\usepackage{booktabs}
\usepackage{hyperref}
\usepackage{enumitem}
\geometry{margin=1in}

\title{\textbf{Our Corrections Analysis}\\
\large Response to MNRAS Editorial Decision MN-26-0837-L}
\author{Eric D. Martin \\ ORCID: 0009-0006-5944-1742}
\date{2026-04-01}

\begin{document}
\maketitle

This document maps each editor objection to the specific change made
in the resubmission (Document 4). Read the rejection letter
(Document 2) alongside this document.

\bigskip
\hrule
\bigskip

\section*{Issue 1 — ``The use of the horizon-normalised coordinate is
simply a redefinition after which no properly-calculated observable
result could change.''}

\textbf{Our analysis:} This is the central objection and reflects a
misreading of the claim. We do \emph{not} claim that a coordinate
transformation changes a physical observable. We claim that the
standard comparison procedure treats two best-fit parameters from
\emph{different cosmological models} as if they were independent
measurements of the same parameter. They are not: $H_0^{\rm SH0ES}$
is extracted within a framework with $\Omega_m = 0.334$;
$H_0^{\rm Planck}$ within $\Omega_m = 0.315$. The 5$\sigma$ tension
conflates this model-space difference with genuine physics.
$\xi = d_c/d_H$ is not a redefinition --- it is a projection onto
the subspace of distance space in which $H_0$ is genuinely absent,
proven analytically by cancellation. What remains in $\xi$ space is
only the $\Omega_m$ residual.

\textbf{Where addressed in Document 4:}
Introduction, paragraph 3 (explicit statement of what the claim is
and is not); Section 2 in full (mathematical demonstration that the
incompatibility is in the comparison procedure, not the coordinates);
Equations 1--6 (showing the exact cancellation and what remains).

\bigskip
\hrule
\bigskip

\section*{Issue 2 — ``Does not point to exactly where [the errors]
are (ideally, specific equations in specific key papers).''}

\textbf{Our analysis:} The original submission discussed the
frame-mixing problem in prose without pointing to specific equations
in the cited papers. This is a legitimate methodological criticism.

\textbf{Where addressed in Document 4:}
Section 2 explicitly cites Riess et al.\ (2022), their Section~2,
for the distance modulus equation (Eq.\ 1 and 2 in Document 4) and
Planck Collaboration (2020), their Section~5.4 and Table~2, for the
acoustic scale relation (Eq.\ 3 in Document 4). These show precisely
where the $H_0$/$\Omega_m$ coupling enters each pipeline.

A footnote in Section 2 clarifies the distinction between the SH0ES
Hubble flow prior ($\Omega_m = 0.30$, fixed in Riess 2022) and the
full Pantheon+ combined best-fit ($\Omega_m = 0.334$, from Brout et
al.\ 2022), which is the value used in this analysis. The frame-mixing
argument holds for any non-zero $\Omega_m$ difference; the quantitative
artifact fraction changes by $<0.5$\% between the two values.

\bigskip
\hrule
\bigskip

\section*{Issue 3 — ``No indication of how [results] were computed,
nor why the Cepheid uncertainty has become much larger.''}

\textbf{Our analysis:} The original Table 1 presented numbers without
derivation. The Cepheid uncertainty of $\pm13.97$ km/s/Mpc compared
to SH0ES's $\pm1.04$ km/s/Mpc required explanation.

\textbf{Where addressed in Document 4:}
Section 4 (new, ``Data and Statistical Methodology'') provides the
complete computation:
\begin{itemize}[noitemsep]
    \item Adaptive quadrature formula for $d_c$ (Eq.~7 in Document 4)
    \item Artifact fraction formula with derivation (Eq.~9)
    \item Cepheid per-object $H_0$ computation (Eq.~10)
\end{itemize}
Section 4.3 explicitly states: the $\pm13.97$ km/s/Mpc is the
\emph{sample standard deviation} across 43 independently analyzed
calibrators. The SH0ES $\pm1.04$ km/s/Mpc is the uncertainty on a
\emph{global pipeline parameter} from a simultaneous fit of 1,701
objects. These are fundamentally different quantities and cannot be
compared directly. The larger value is the statistically honest one
for individual-calibrator inference.

Every row in the revised Table 1 references the equation that
produced it.

\bigskip
\hrule
\bigskip

\section*{Issue 4 — ``Nor is there indication of the statistical
methodology being used.''}

\textbf{Our analysis:} No explicit statistical methodology section
existed in the original submission.

\textbf{Where addressed in Document 4:}
Section 4 provides:
\begin{itemize}[noitemsep]
    \item Dataset description (1,701 SNe~Ia, 43 Cepheid calibrators,
    public Pantheon+SH0ES release)
    \item Computation method (adaptive quadrature, two cosmologies)
    \item Artifact fraction formula (Eq.~9) with bootstrap uncertainty
    quantification ($10^4$ iterations, $\pm0.1$\% standard error)
    \item Explicit Cepheid analysis pipeline (Eq.~10)
\end{itemize}

\bigskip
\hrule
\bigskip

\section*{Issue 5 — ``Overall the article is not constructed or
presented at a level suitable for a professional international
astronomy journal.''}

\textbf{Our analysis:} Two components: (a) formatting, (b) structure.

\textbf{Where addressed in Document 4:}
\begin{itemize}[noitemsep]
    \item Document class changed from generic \texttt{article} to
    \texttt{mnras} (official MNRAS LaTeX class)
    \item Session identifiers removed from the date line
    \item New Section 2 added, providing the standard ``problem
    statement'' that was missing from the original
    \item New Section 4 (methodology) elevates the paper to standard
    observational cosmology paper structure:
    Introduction / Problem Statement / Method / Data \& Statistics /
    Results / Discussion / Reply to Objections / Conclusion
    \item All five analysis scripts referenced by name with public
    repository URL for full reproducibility
\end{itemize}

\bigskip
\hrule
\bigskip

\section*{What Was Not Changed}

The following were deliberately preserved:
\begin{itemize}[noitemsep]
    \item All numerical results (93.3\%, 3\% residual, Cepheid
    $H_0 \approx 70$) --- these are correct and confirmed
    \item The JWST, TRGB, and EDE reply sections --- these directly
    address the strongest known objections to the claim
    \item The DESI BAO convergent $\Omega_m$ evidence --- this is
    independent confirmation that the 3\% residual is real
    \item The Forward Implication section --- this is the practical
    argument for why the field should adopt $\xi$ normalization
\end{itemize}

\section*{The Single Remaining Risk}

A reviewer may still argue: ``Even if the comparison procedure is
imperfect, the tension has been confirmed by multiple independent
pipelines (TRGB, masers, lensing) that do not share the SH0ES
coordinate assumptions.'' This objection is addressed in Section 6.2
of Document 4, which notes that the TRGB results are not uniform
(Freedman JWST TRGB gives $H_0 \approx 70$) and that the systematic
bifurcation tracks pipeline membership, not fundamental physics.
A reviewer who pushes on this point is engaging with the science ---
that is the correct outcome.

\end{document}
